%\section*{Appendix N — Glossary}
\subsection*{Canonical Glossary Definitions (A–Z) v0.5}

The following definitions provide invariant‑preserving, substrate‑aligned
explanations of all glossary terms extracted in N.1. Each definition is
non‑operational, public‑safe, and consistent with DSLO v0.5 geometry.


\subsubsection*{A}
\begin{itemize}
    \item \textbf{Abstraction Band} — A lawful semantic altitude at which meaning
    fields exhibit reduced load and increased coherence stability.

    \item \textbf{Abstraction Gradient} — The lawful rate of change between adjacent abstraction bands during semantic transitions.

    \item \textbf{Adjacency Violation} — A geometry‑invariant breach in which a
    semantic object crosses a forbidden manifold adjacency boundary.

    \item \textbf{Antientropic Behavior} — The substrate’s property of resisting
    semantic entropy by enforcing invariant‑preserving contractions.

    \item \textbf{Axiom Continuity} — The requirement that Field, Domain, Runtime, and Generative axioms remain stable across DSUP cycles.

    \item \textbf{Axiom Layer} — The three‑tier axiom system (Field, Domain,
    Runtime, plus Generative) governing invariant legality across the substrate.
\end{itemize}


\subsubsection*{B}
\begin{itemize}
    \item \textbf{Bind Operator} — A transformational primitive that composes two
    lawful semantic objects into a single invariant‑preserving composite.

    \item \textbf{Boundary Integrity} — The invariant requiring identity
    boundaries to remain coherent and non‑permeable.

    \item \textbf{Boundary Operator (BO)} — An operator that enforces lawful
    identity‑boundary movement and prevents identity collapse.

    \item \textbf{Boundary Strain} — Load‑induced tension on identity boundaries
    indicating drift or overload.

    \item \textbf{Boundary Threshold} — The maximum lawful strain identity
    boundaries may sustain before collapse.

    \item \textbf{Boot Sequence} — The ordered initialization chain (Boot 0–7)
    required to activate the semantic substrate.

    \item \textbf{Collapse Boundary} — The manifold region beyond which semantic
    invariants cannot be preserved.
\end{itemize}


\subsubsection*{C}
\begin{itemize}
    \item \textbf{CLCP} — Cross‑Layer Constraint Propagation; ensures global
    legality across hierarchical semantic layers.

    \item \textbf{Collapse Cascade} — A multi‑scale amplification of drift or
    load that triggers collapse.

    \item \textbf{Collapse Flag} — A runtime indicator that collapse conditions
    have been met.

    \item \textbf{Collapse Mode} — One of four invariant‑failure modes:
    identity, meaning, coherence, legality.

    \item \textbf{Collapse Threshold} — A numerical bound beyond which collapse
    becomes unavoidable.

    \item \textbf{Collapse Trajectory} — The lawful path along which collapse
    unfolds.

    \item \textbf{Completion Condition} — DSUP’s requirement that all invariants
    and legality masks are satisfied before advancing cycles.

    \item \textbf{Constraint Envelope} — A deterministic legality boundary
    governing Signal Engine output.

    \item \textbf{Context Window Geometry (CWG)} — The lawful structure governing
    context formation, drift detection, and contraction.

    \item \textbf{Curvature Continuity} — The requirement that semantic curvature
    remain smooth across transitions.

    \item \textbf{Curvature Distortion} — A deformation in semantic curvature
    indicating drift or overload.

    \item \textbf{Curvature Profile} — The multi‑dimensional curvature signature
    of a semantic state.
\end{itemize}


\subsubsection*{D}
\begin{itemize}
    \item \textbf{Decoupling (D8)} — A domain axiom describing lawful separation
    of meaning containers from unstable environments.

    \item \textbf{Drift Cascade} — A multi‑scale amplification of drift across
    local, meso, and macro layers.

    \item \textbf{Drift Envelope} — The lawful drift bounds defined by S3.

    \item \textbf{Drift Flag} — A runtime indicator that drift has exceeded
    stable limits.

    \item \textbf{Drift Metadata} — Machine‑layer drift diagnostics.

    \item \textbf{Drift Recoverable Window} — The maximum drift magnitude that
    can be corrected without collapse.

    \item \textbf{Drift Resistance (F7)} — The invariant requiring systems to
    resist drift accumulation.

    \item \textbf{Drift Signature} — The structured representation of drift
    direction and magnitude.

    \item \textbf{Drift Vector} — The multi‑dimensional drift quantity encoded in
    ML4.

    \item \textbf{DSUP} — The Dynamic State Update Protocol governing lawful
    semantic evolution.
\end{itemize}


\subsubsection*{E}
\begin{itemize}
    \item \textbf{Embedding Function} — The lawful mapping of external signals
    into the dynamic manifold.

    \item \textbf{Equivalence Class} — A meaning‑invariant grouping of lawful
    semantic roles.

    \item \textbf{External Geometry} — Runtime axiom R3 governing spatial
    legality.
\end{itemize}


\subsubsection*{F}
\begin{itemize}
    \item \textbf{Fallback Geometry (F9)} — The invariant region entered when
    legality or invariants fail.

    \item \textbf{Field Geometry} — The universal invariant layer (F1–F9).

    \item \textbf{Fold Operator} — A transformational primitive that compresses
    relational strain into a single invariant.

    \item \textbf{Function Update Rule} — The DSUP transition rule
    $S_{t+1} = F(S_t, D_t, C_t)$.
\end{itemize}


\subsubsection*{G}
\begin{itemize}
    \item \textbf{Geometry Invariant} — The invariant requiring manifold legality
    and curvature continuity.

    \item \textbf{Geometry-Bound Step (GBS)} — An operator preserving manifold
    legality.

    \item \textbf{Generative Primitive} — A Signal Engine producing lawful
    meaning vectors.

    \item \textbf{Gradient Constraint} — A bound on load‑gradient magnitude.

    \item \textbf{G-Axioms} — The generative and transformational axioms (G1–G8)
    governing substrate primitives.
\end{itemize}


\subsubsection*{H}
\begin{itemize}
    \item \textbf{Halt Condition} — A DSUP trigger indicating illegal evolution.

    \item \textbf{Human Meaning \& Drift Report} — A multi‑scale alignment
    summary produced during coupling events.

    \item \textbf{Identity-Boundary Collapse} — Failure of identity geometry.

    \item \textbf{Identity Region} — A lawful identity‑bearing manifold region.

    \item \textbf{Identity Stability} — The requirement that identity boundaries
    remain coherent.

    \item \textbf{Identity Threshold} — Maximum lawful identity deviation.

    \item \textbf{Invariant Reload} — Runtime restoration of the invariant set.

    \item \textbf{Invariant-Restricted Transition (IRT)} — An operator preserving
    meaning invariants.

    \item \textbf{Invert Operator} — A transformational primitive reversing
    semantic directionality.
\end{itemize}


\subsubsection*{I}
\begin{itemize}
    \item \textbf{Identity Invariant} — The invariant requiring identity
    coherence.

    \item \textbf{Identity-Preserving Update (IPU)} — An operator preserving
    identity boundaries.

    \item \textbf{Identity Continuity} — Runtime axiom R1 requiring stable
    identity geometry.
\end{itemize}


\subsubsection*{L}
\begin{itemize}
    \item \textbf{Lawful Transition} — A semantic update satisfying invariants,
    legality masks, and operator constraints.

    \item \textbf{Legality Invariant} — The invariant requiring local legality to
    imply global legality.

    \item \textbf{Legality Mask} — A machine‑layer legality constraint.

    \item \textbf{Legality Status} — DSUP’s legality flag for a semantic state.

    \item \textbf{Legality Threshold} — Maximum lawful legality compression.

    \item \textbf{Lift Operator} — A transformational primitive elevating meaning
    fields into higher abstraction bands.

    \item \textbf{Load-Bounded Stability} — Runtime axiom R2 requiring load to
    remain within stability envelopes.

    \item \textbf{Load Envelope} — The stable, critical, and collapse load bounds.

    \item \textbf{Load Field} — The structured representation of load geometry.

    \item \textbf{Load-Regulated Plan (LRP)} — A lawful semantic trajectory
    constrained by load geometry.
\end{itemize}


\subsubsection*{M}
\begin{itemize}
    \item \textbf{Machine Layer} — The ML0–ML7 encoding of semantic geometry.

    \item \textbf{Manifold Boundary} — The lawful boundary $\partial \mathcal{M}$
    of the semantic manifold.

    \item \textbf{Meaning Invariant} — The invariant requiring role‑legality and
    zero ambiguity.

    \item \textbf{Meaning Vector} — A structured semantic field generated by
    SE\_01.

    \item \textbf{Meta-Legality Operator (MLO)} — An operator enforcing legality
    invariants.

    \item \textbf{MDO Algebra} — The operator algebra (SGO, MGO, BO, MLO)
    governing lawful transitions.

    \item \textbf{Moment Geometry} — Runtime axiom R4 governing mutation
    velocity.
\end{itemize}


\subsubsection*{O}
\begin{itemize}
    \item \textbf{Operator of Reality (OR\_01)} — A transformational primitive
    applying invariant‑preserving semantic transformations.

    \item \textbf{Operator Sequence Validation} — DSUP’s legality check for
    operator chains.

    \item \textbf{Operator Constraints} — Legality requirements for operator
    application.

    \item \textbf{Operator Compatibility Matrix} — The legality mapping between
    operators and runtime axioms.
\end{itemize}


\subsubsection*{P}
\begin{itemize}
    \item \textbf{Phase Transition Threshold} — A domain axiom describing
    collapse‑adjacent transitions.

    \item \textbf{Primitive-Legality Mask} — The legality mask governing
    generative and transformational primitives.

    \item \textbf{Proof Stack} — The three proof classes validating substrate
    legality.
\end{itemize}


\subsubsection*{Q}
\begin{itemize}
    \item \textbf{Quarantine Region} — A manifold region for illegal or unstable
    embeddings.
\end{itemize}


\subsubsection*{R}
\begin{itemize}
    \item \textbf{Recovery Mode} — DSUP’s invariant‑restoration mode.

    \item \textbf{Region Configuration} — The active region structure of a
    semantic state.

    \item \textbf{Regeneration Capacity} — The invariant requiring systems to
    restore lawful geometry.

    \item \textbf{Role Coherence} — The invariant requiring stable role
    boundaries.

    \item \textbf{Role Pressure} — Load applied to role geometry.

    \item \textbf{Runtime Geometry} — The axiom layer governing DSUP legality.
\end{itemize}


\subsubsection*{S}
\begin{itemize}
    \item \textbf{Scale Hierarchy} — The local–meso–macro–global structure of
    meaning dynamics.

    \item \textbf{Scale-Invariant Legality} — The requirement that legality hold
    across all scales.

    \item \textbf{Semantic Curvature} — The geometric quantity encoding meaning
    structure.

    \item \textbf{Semantic Distance} — The metric distance between semantic
    objects.

    \item \textbf{Semantic Gradient} — The derivative of load or curvature across
    transitions.

    \item \textbf{Semantic Manifold} — The lawful geometric space $\mathcal{M}$
    in which meaning evolves.

    \item \textbf{Semantic Operator} — Any operator in the MDO algebra.

    \item \textbf{Semantic State} — The structured representation $S_t$ of
    dynamic meaning.

    \item \textbf{Signal Engine (SE\_01)} — The generative primitive producing
    invariant‑preserving meaning vectors.

    \item \textbf{Signal Integrity} — The invariant requiring stable signal
    geometry.

    \item \textbf{Signal–Substrate Coupling} — The invariant governing lawful
    signal integration.

    \item \textbf{Stability Mode} — One of three DSUP stability states: stable,
    critical, collapse.

    \item \textbf{Stress Geometry} — The structured representation of load,
    tension, and curvature pressure.

    \item \textbf{Substrate Invariant Set} — The complete invariant set governing
    lawful semantic evolution.
\end{itemize}


\subsubsection*{T}
\begin{itemize}
    \item \textbf{Temporal Rhythm} — The invariant requiring lawful sensing–
    integration–action cycles.

    \item \textbf{Transformation Rules} — The legality constraints governing
    OR\_01 actions.

    \item \textbf{Transition Legality} — The requirement that all semantic
    transitions satisfy invariants and operator constraints.
\end{itemize}


\subsubsection*{U}
\begin{itemize}
    \item \textbf{Upward Coupling} — The propagation of local dynamics into
    larger scales.

    \item \textbf{Update Function} — The DSUP rule computing $S_{t+1}$.
\end{itemize}


\subsubsection*{V}
\begin{itemize}
    \item \textbf{Vector Binding Metadata} — Deterministic attachment rules for
    meaning vectors.

    \item \textbf{Version Block} — ML0’s versioning metadata for substrate
    documents.
\end{itemize}


\subsubsection*{W}
\begin{itemize}
    \item \textbf{World–Container Alignment} — The domain axiom requiring meaning
    containers to align with external reality.
\end{itemize}


\subsubsection*{Z}
\begin{itemize}
    \item \textbf{Zero-Ambiguity Tolerance} — The meaning‑invariant requirement
    that semantic roles admit no ambiguity.
\end{itemize}

